\par\bigskip
\hypertarget{native-395}{}
\begin{center}\pdfbookmark[2]{6. Secularstörungen}{ceres-006}[6.]\end{center}
\noindent\hspace*{\parindent}[Man kann nun alle Glieder bilden, aus denen sich \(\delta r\) und \(\delta v\) zusammensetzen, indem man in die Gleichung 14 die einzelnen Theile von \(Q\) nach 13 einsetzt.] So hat man die Theile von \(\delta r\) folgendermassen
\[
\begin{gathered}
1)\quad\text{aus }2g+\mathfrak A^{(0)}\quad .\, .\, .\, .\, .\, .\, .\, .\, .\, .\, .\, .\\[.5ex]
-aa(2g+\mathfrak A^{(0)})\{1-\tfrac14e\cos M+\underline{\tfrac32ent\sin M\}}
\end{gathered}
\]
\[
\begin{gathered}
2)\quad\text{aus }\mathfrak A^{(i)}\cos iD\quad .\, .\, .\, .\, .\, .\, .\, .\, .\, .\, .\, .\, .\, .\\[.5ex]
aa\mathfrak A^{(i)}\left\{\begin{aligned}
&\frac1{hh-1}\cos iD\\
&+\tfrac14e\left(\frac1{h+2}-\frac{4h}{hh-1}+\frac3h\right)\cos(iD+M)\\
&-\tfrac14e\left(\frac1{h-2}-\frac{4h}{hh-1}+\frac3h\right)\cos(iD-M)
\end{aligned}\right\}
\end{gathered}
\]
\[
\begin{gathered}
3)\quad\text{aus }e\mathfrak B^{(0)}\cos M\quad .\, .\, .\, .\, .\, .\, .\, .\, .\, .\, .\\[.5ex]
-eaa\mathfrak B^{(0)}\{\tfrac14\cos M+\underline{\tfrac12nt\sin M}\}
\end{gathered}
\]
\[
\begin{gathered}
4)\quad\text{aus }e'\mathfrak C^{(0)}\cos(D+M')\quad .\, .\, .\, .\, .\\[.5ex]
-e'aa\mathfrak C^{(0)}\{\tfrac14\cos(D+M')+\underline{\tfrac12nt\sin(D+M')}\}
\end{gathered}
\]
\[
\begin{gathered}
5)\quad\text{aus }e\mathfrak B^{(i)}\cos(iD+M)\quad .\, .\, .\, .\, .\, .\, .\\[.5ex]
eaa\mathfrak B^{(i)}\cdot\frac1{hh+2h}\cos(iD+M)
\end{gathered}
\]
\[
\begin{gathered}
6)\quad\text{aus }e'\mathfrak C^{(i)}\cos((i+1)D+M')\quad .\, .\\[.5ex]
e'aa\mathfrak C^{(i)}\cdot\frac1{hh+2h}\cos((i+1)D+M').
\end{gathered}
\]

[Die Summe aller dieser Theile, welche noch mit der störenden Masse \(\mu\) zu multipliciren sind, ist \(\delta r\), wobei in 2, 5, 6 \(i\) alle positiven und negativen ganzen Zahlen, Null ausgeschlossen, darstellt.

Hingegen [ergeben sich] die Theile von \(\delta v\) [welche aus dem ersten Gliede der Gleichung 15 stammen, folgendermassen:]
\[
\begin{gathered}
1)\quad\text{aus }2g+A^{(0)}\quad .\, .\, .\, .\, .\, .\, .\, .\, .\, .\, .\, .\\[.5ex]
-a(2g+\mathfrak A^{(0)})\{\tfrac52e\sin M+\underline{3ent\cos M\}}
\end{gathered}
\]
\[
\begin{gathered}
2)\quad\text{aus }\mathfrak A^{(i)}\cos iD\quad .\, .\, .\, .\, .\, .\, .\, .\, .\, .\, .\, .\\[.5ex]
-a\mathfrak A^{(i)}\left\{\begin{aligned}
&\frac{2h}{hh-1}\sin iD\\
&+\left(\frac{\tfrac12(h+1)}{h+2}-\frac{2hh+h-\tfrac12}{hh-1}+\frac{\tfrac32(h+1)}h\right)e\sin(iD+M)\\
&-\left(\frac{\tfrac12(h-1)}{h-2}-\frac{2hh-h-\tfrac12}{hh-1}+\frac{\tfrac32(h-1)}h\right)e\sin(iD-M)
\end{aligned}\right\}
\end{gathered}
\]
\[
\begin{gathered}
3)\quad\text{aus }e\mathfrak B^{(0)}\cos M\quad .\, .\, .\, .\, .\, .\, .\, .\, .\, .\, .\\[.5ex]
-\tfrac12ea\mathfrak B^{(0)}\{\sin M+\underline{2nt\cos M}\}
\end{gathered}
\]
\[
\begin{gathered}
4)\quad\text{aus }e'\mathfrak C^{(0)}\cos(D+M')\quad .\, .\, .\, .\, .\\[.5ex]
-\tfrac12ea\mathfrak C^{(0)}\{\sin(D+M')+\underline{2nt\cos(D+M')}\}
\end{gathered}
\]
\[
\begin{gathered}
5)\quad\text{aus }e\mathfrak B^{(i)}\cos(iD+M)\quad .\, .\, .\, .\, .\\[.5ex]
-2ea\mathfrak B^{(i)}\frac{h+1}{hh+2h}\sin(iD+M)
\end{gathered}
\]
\[
\begin{gathered}
6)\quad\text{aus }e'\mathfrak C^{(i)}\cos((i+1)D+M')\quad .\, .\\[.5ex]
-2e'a\mathfrak C^{(i)}\frac{h+1}{hh+2h}\sin((i+1)D+M').
\end{gathered}
\]

[Die aus den beiden letzten Gliedern von 15] noch hinzuzusetzenden Theile sind:
\hypertarget{native-396}{}
\[
\begin{aligned}
&a(3g+2\mathfrak A^{(0)})nt+\Sigma\frac{a(2\mathfrak A^{(i)}-\tfrac12\alpha^{(i)})}h\sin iD\\
&+(2\mathfrak B^{(0)}-\tfrac12\beta^{(0)})ea\sin M+\Sigma\frac{2\mathfrak B^{(i)}-\tfrac12\beta^{(i)}}{h+1}ea\sin(iD+M)\\
&+(2\mathfrak C^{(0)}-\tfrac12\gamma^{(0)})e'a\sin(D+M')\\
&\qquad+\Sigma\frac{2\mathfrak C^{(i)}-\tfrac12\gamma^{(i)}}{h+1}e'a\sin((i+1)D+M'),
\end{aligned}
\]
\noindent[wo der Coefficient \(\alpha^{(0)}\) bei Seite gelassen ist, weil für den constanten Theil von \(\int\partial R\) die Bezeichnung \(g\) eingeführt ist. Auch hier sind alle Glieder noch mit \(\mu\) zu multipliciren.]

[Es sollen nun noch einige passende Transformationen vorgenommen werden:

In \(\delta v\) ist das Glied \(a(3g+2\mathfrak A^{(0)})nt\) gefunden worden, das durch passende Wahl von \(g\) zum Verschwinden gebracht werden kann, so dass \(n\) in der That die Bedeutung der mittlern Bewegung erlangt;] hieraus folgt
\[g=-\tfrac13a\left(\frac{\mathrm{d}A^{(0)}}{\mathrm{d}a}\right)\mathpunct{\cdot}\]

[Weiter sollen die secularen Glieder in \(\delta r\) und \(\delta v\) in Secularstörungen der Excentricität und der Sonnenferne verwandelt werden; es ist nach den obigen Formeln
\[
\begin{gathered}
\text{der seculare Theil von }\delta r\text{ gleich}\\[.5ex]
\tfrac14aa(B^{(0)}-2\mathfrak B^{(0)})ent\sin M-\tfrac12aa\mathfrak C^{(0)}e'nt\sin(D+M')\\[1ex]
\text{»}\qquad\text{»}\qquad\text{»}\qquad\text{»}\quad\delta v\quad\text{»}\\[.5ex]
\tfrac12a(\mathfrak B^{(0)}-2\mathfrak B^{(0)})ent\cos M-a\mathfrak C^{(0)}e'nt\cos(D+M')
\end{gathered}
\]
\noindent und, da \(D=M-M'+\tilde\omega-\tilde\omega'\), so wird
\[
\begin{gathered}
\text{der erstere gleich}\\[.5ex]
\begin{aligned}
&\tfrac14aa\left(B^{(0)}-2\mathfrak B^{(0)}-2\frac{e'}e\mathfrak C^{(0)}\cos(\tilde\omega-\tilde\omega')\right)ent\sin M\\
&\hfill-\tfrac12aa\mathfrak C^{(0)}\sin(\tilde\omega-\tilde\omega')e'nt\cos M
\end{aligned}
\end{gathered}
\]
\[
\begin{gathered}
\text{der letztere gleich}\\[.5ex]
\begin{aligned}
&\tfrac12a\left(B^{(0)}-2\mathfrak B^{(0)}-2\frac{e'}e\mathfrak C^{(0)}\cos(\tilde\omega-\tilde\omega')\right)ent\cos M\\
&\hfill+a\mathfrak C^{(0)}\sin(\tilde\omega-\tilde\omega')e'nt\sin M.
\end{aligned}
\end{gathered}
\]

Aus den Relationen
\[r=a(1+e\cos M)\qquad v=nt+\epsilon-2e\sin M\]
\noindent erhält man aber
\[
\begin{aligned}
\delta r&=a\cos M\mathord{.}\,\delta e+ae\sin M\mathord{.}\,\delta\tilde\omega\\
\delta v&=-2\sin M\mathord{.}\,\delta e+2e\cos M\mathord{.}\,\delta\tilde\omega;
\end{aligned}
\]
\noindent\hypertarget{native-397}{}also wird
\[
\begin{aligned}
\delta e&=\tfrac12e'a\mathfrak C^{(0)}nt\sin(\tilde\omega'-\tilde\omega)\\
\delta\tilde\omega&=\tfrac14a\left(B^{(0)}-2\mathfrak B^{(0)}-2\frac{e'}e\mathfrak C^{(0)}\cos(\tilde\omega'-\tilde\omega)\right)nt.]
\end{aligned}
\]
\par\medskip\centerline{\rule{.28\linewidth}{.35pt}}\par
